\documentclass[11pt,a4paper]{article}
\usepackage{fontspec}
\setmainfont{Latin Modern Roman}
\usepackage[margin=23mm]{geometry}
\usepackage{amsmath,amssymb,amsthm,mathtools,mathrsfs}
\usepackage{longtable,booktabs}
\usepackage[unicode,hidelinks]{hyperref}
\providecommand{\C}{\mathbb C}
\providecommand{\R}{\mathbb R}
\newtheorem{theorem}{Theorem}[section]
\newtheorem{lemma}[theorem]{Lemma}
\newtheorem{proposition}[theorem]{Proposition}
\newtheorem{corollary}[theorem]{Corollary}
\newtheorem{definition}[theorem]{Definition}
\newtheorem{remark}[theorem]{Remark}
\allowdisplaybreaks
\emergencystretch=3em
\makeatletter
\renewcommand\@pnumwidth{2.5em}
\renewcommand\@tocrmarg{3.5em}
\makeatother
\title{The Literal Cyclic Sectors and Original Mixed Kernel}
\author{Split-Zero research programme}
\date{14 September 2026}
\begin{document}
\maketitle
\paragraph{Portable source edition.}
Every TeX input used by this entry is expanded below. Historical file
and directory names in the preserved text are source locators; the
exact origin paths and hashes are recorded in \texttt{11\_SOURCE\_MAP.json}.
The other numbered proof files in this flat handoff supply the
accompanying complete calculations.\par

% New original-observation calculation; the sealed cut22 is not edited.
\section{The literal cyclic sectors and the original mixed kernel}
\label{sec:OCS}

This calculation evaluates the cyclic projector already present in
the original period observation. Its coefficients retain the original
period integrals, amplitude unit, and all primary nilpotents. The
result includes an exact finite matrix for the actual rank, as well
as unconditional rank bounds obtained from the original symmetric
tensor geometry.

\subsection{Original data and the Fourier basis of the literal matrix}
Use the following data of OPG1--5, including the same contours and
their orientations:
\[
 \begin{gathered}
 \alpha_{\epsilon,\eta}=\tfrac12+\epsilon\delta+i\eta\gamma,
 \quad\epsilon,\eta\in\{+1,-1\},\quad0<\delta<\tfrac12,\quad\gamma>2,\\
 h(z_1)=\prod_{\epsilon,\eta}(z_1-\alpha_{\epsilon,\eta})^m,
 \quad n=4m,\quad N=n+1,\quad E_h=\mathbb C[z_1]/(h),\\
 k=4l+1,\quad l,m\geq1,\quad c=k/2,
 \quad e=1+k(m-1),\quad q=e(k+1)^2,\\
 \beta_{a,b}=c+(2a-k)\delta+i(2b-k)\gamma,
 \quad\chi(S)=\prod_{a,b=0}^k(S-\beta_{a,b})^e,
 \quad E=\mathbb C[S]/(\chi).
 \end{gathered}\tag{OCS1}
\]
Here $A_h$ and $A$ are multiplication by $z_1$ and $S$ on their
respective algebras, and $Y=(A-cI)/i$. The original amplitude
unit is $U=M_{j_hv_h}$, with $v_h=(2\xi)/h$ and every original
$v_h(\alpha)\ne0$. For the specified nonzero $u$ and its branch,
$\Phi_h'=h$, $\Phi_h(0)=0$, and
\[
 \Gamma_j=-\ell_0+\ell_j,\qquad
 \arg\ell_j=\frac{\arg u+\pi+2\pi j}{N},\qquad
 (\Pi[P])_j=\int_{\Gamma_j}e^{\Phi_h(z_1)/u}
                       \operatorname{rem}_hP(z_1)\,dz_1.
 \tag{OCS2}
\]
The representative has degree less than $n$. The convergence and
invertibility of this original period map are the complete period
theorems retained in OPG4. No period value or unit derivative is
chosen in the present calculation. The original injection and
observation are
\[
 \iota[P]=[P(z_1+\cdots+z_k)],\quad
 \mathcal V=\Pi^{\otimes k}U^{\otimes k}\iota,
 \quad P_k=\frac1N\sum_{a=0}^{N-1}(C^{-a})^{\otimes k},
 \quad L=P_k\mathcal V=\jmath\Lambda,
 \quad B=\operatorname{im}L,\quad K=\ker\Lambda.
 \tag{OCS3}
\]
The map $\jmath:B\hookrightarrow(\mathbb C^n)^{\otimes k}$
is the original image inclusion. The injection $\iota$ has precisely
the defining polynomial in OCS1: on each ordered primary tensor
the sum of its $k$ independent nilpotents has index
$e=1+k(m-1)$, its top coefficient being
$(k(m-1))!/((m-1)!)^k\ne0$, and its distinct eigenvalues are
the displayed $\beta_{a,b}$. Thus $\mathcal V$ is injective.

On a period-coordinate column $x\in\mathbb C^{N-1}$ the literal
matrix in OPG4--5 is
\[
 (Cx)_j=x_{j+1}-x_1\quad(1\leq j<N-1),\qquad
 (Cx)_{N-1}=-x_1.
 \tag{OCS4}
\]
Set $\zeta=\exp(2\pi i/N)$ and, for $1\leq t\leq N-1$, let
\[
 (f_t)_j=\zeta^{jt}-1,\qquad
 F=[f_1\ \cdots\ f_{N-1}],\qquad
 (F^{-1}x)_t=\frac1N\sum_{j=1}^{N-1}\zeta^{-tj}x_j.
 \tag{OCS5}
\]
The last expression is the exact inverse. Indeed
$\sum_{j=1}^{N-1}\zeta^{-tj}(\zeta^{uj}-1)=N\mathbf1_{t=u}$
for $t,u\ne0$ modulo $N$, by the finite geometric sum.
The first $N-2$ rows of OCS4 give
$\zeta^{(j+1)t}-\zeta^t=\zeta^t(\zeta^{jt}-1)$, and the
last row gives $1-\zeta^t=\zeta^t(\zeta^{(N-1)t}-1)$.
Consequently $Cf_t=\zeta^tf_t$, proving the orientation of
every inverse power in OCS3, as well as $C^N=I$.
The original Euclidean Gram of this basis is
\[
 F^*F=N(I_{N-1}+\mathbf1\mathbf1^*).
 \tag{OCS6}
\]
To verify it, expand
$\sum_j(\zeta^{-jt}-1)(\zeta^{ju}-1)$ and use the same
geometric sums; its value is $N(\mathbf1_{t=u}+1)$.
Thus the complete coefficient and metric transformations are
both specified.

\subsection{The ordered-to-symmetric map and the exact sector projector}
Let
\[
 \mathcal A_k=\{\nu\in\mathbb Z_{\geq0}^{N-1}:\sum_t\nu_t=k\},
 \quad\operatorname{wt}(\nu)=\sum_{t=1}^{N-1}t\nu_t\pmod N,
 \quad\mathcal I_{k,\tau}=\{\nu\in\mathcal A_k:
                                      \operatorname{wt}(\nu)=\tau\}.
 \tag{OCS7}
\]
For each $\nu$ define the following unscaled symmetric tensor:
\[
 S_\nu=\sum_{\substack{(t_1,\ldots,t_k)\in\{1,\ldots,N-1\}^k\\
                   \#\{j:t_j=t\}=\nu_t\ \text{for each }t}}
                    f_{t_1}\otimes\cdots\otimes f_{t_k}.
 \tag{OCS8}
\]
Every distinct ordered word occurs once. Since the ordered tensors
in the $f_t$ form a basis, their disjoint word supports show that
the $S_\nu$ are linearly independent. A permutation-invariant tensor
has equal coefficients on each word orbit, so they span exactly
$\operatorname{Sym}^k(\mathbb C^{N-1})$, realized as that invariant
subspace of the full ordered tensor product. In particular, for
$x=(x_1,\ldots,x_{N-1})$,
\[
 (Fx)^{\otimes k}=\sum_{\nu\in\mathcal A_k}x^\nu S_\nu,
 \quad
 (C^{-a})^{\otimes k}S_\nu=
                         \zeta^{-a\operatorname{wt}(\nu)}S_\nu,
 \quad
 P_k(Fx)^{\otimes k}=\sum_{\nu\in\mathcal I_{k,0}}x^\nu S_\nu.
 \tag{OCS9}
\]
The finite average in OCS3 is exactly one on weight zero and
zero on every other weight, by its geometric sum. No multinomial
coefficient is omitted: the multiplicity $k!/\prod_t\nu_t!$
is already the number of summands in OCS8.

Write $\mathcal F_k:\mathbb C^{\mathcal A_k}\to
\operatorname{Sym}^k(\mathbb C^{N-1})$ for the basis isomorphism
$e_\nu\mapsto S_\nu$. Let $J_0$ be the coordinate inclusion of
$\mathbb C^{\mathcal I_{k,0}}$ into $\mathbb C^{\mathcal A_k}$,
$\pi_0$ the coordinate restriction in the opposite direction,
and let $\mathcal F_0=\mathcal F_kJ_0$, with its target the original
ordered tensor space. The exact projector factorization is
\[
 P_k\mathcal F_k=\mathcal F_0\pi_0.
 \tag{OCS10}
\]
All images of $\mathcal V$ are symmetric: polynomial evaluation
at a sum is invariant under every factor permutation, and the
same map $\Pi U$ is applied in all factors. Therefore OCS10
applies on the entire original image $\mathcal V(E)$.

For completeness, the full inherited metric is retained in these
coordinates. Put $G=F^*F$ as in OCS6. The Gram of the basis OCS8 is
\[
 \mathcal G_{\nu,\mu}=
 \sum_{\substack{a_{tu}\geq0\\\sum_u a_{tu}=\nu_t\ ,\
                                    \sum_t a_{tu}=\mu_u}}
       \frac{k!}{\prod_{t,u}a_{tu}!}\prod_{t,u}G_{tu}^{a_{tu}}.
 \tag{OCS11}
\]
Indeed the coefficient of $\bar x^\nu y^\mu$ in
$(\bar x^TGy)^k$ is the right side by the multinomial theorem.
Taking the inner product of the two pure tensors in OCS9 gives
the same coefficient as $\langle S_\nu,S_\mu\rangle$.
This proves OCS11, including its off-diagonal entries; it is
positive definite because OCS8 is a basis. On the full symmetric
coordinates the projector matrix is $J_0\pi_0$, and its adjoint
for this exact Gram is
$\mathcal G^{-1}(J_0\pi_0)^*\mathcal G$.
Thus the Fourier-sector computation makes no unproved
orthogonality or norm assertion about the original cyclic average.

\subsection{Every primary period coefficient of the projected curve}
For each original root $\alpha$ let $E_{\alpha,a}$, $0\leq a<m$,
be the degree-less-than-$4m$ primary representatives of OPG10.
Explicitly, if $h_\alpha=h/(z_1-\alpha)^m$, put
\[
 E_{\alpha,a}(z_1)=\operatorname{rem}_h\left[
 h_\alpha(z_1)\sum_{j=0}^{m-1}
       \frac{(h_\alpha^{-1})^{(j)}(\alpha)}{j!}
                  (z_1-\alpha)^{j+a}\right],\qquad
 b_{j;\alpha,a}=\int_{\Gamma_j}e^{\Phi_h(z_1)/u}
                                    E_{\alpha,a}(z_1)\,dz_1.
 \tag{OCS12}
\]
The four primary factors are coprime. The Taylor inverse gives
constant one on its own factor and zero on the other three;
thus the classes $[E_{\alpha,a}]$ form the full primary basis
with all $m$ nilpotent degrees at every root.
Define the exact finite Fourier and unit coefficients
\[
 \widehat b_{t;\alpha,a}=
     \frac1N\sum_{j=1}^{N-1}\zeta^{-tj}b_{j;\alpha,a},
 \qquad
 f_{t;\alpha,d}=\frac{i^{-d}}{d!}
       \sum_{a=d}^{m-1}\widehat b_{t;\alpha,a}
                         \frac{v_h^{(a-d)}(\alpha)}{(a-d)!},
 \quad0\leq d<m.
 \tag{OCS13}
\]
Let $\omega_\alpha=(\alpha-1/2)/i$ and set
\[
 \mathbf v(w)=\Pi U\exp\bigl(w(A_h-\tfrac12I)/i\bigr)[1],
 \quad x(w)=F^{-1}\mathbf v(w),\qquad
 x_t(w)=\sum_{\alpha}e^{\omega_\alpha w}
                          \sum_{d=0}^{m-1}f_{t;\alpha,d}w^d.
 \tag{OCS14}
\]
To verify the final equality, in the primary algebra at $\alpha$
multiply the Taylor series of $v_h$ by the terminating series of
$\exp(w(z_1-\alpha)/i)$. The coefficient on
$[E_{\alpha,a}]$ is
$e^{\omega_\alpha w}\sum_{d=0}^a
v_h^{(a-d)}(\alpha)(w/i)^d/((a-d)!d!)$.
Apply the original polynomial period map and the inverse in OCS5;
this gives precisely OCS13--14. In particular, the unreduced
entire function has passed through its full primary remainder
before any period is integrated.

For $\nu\in\mathcal I_{k,0}$, $0\leq a,b\leq k$ and
$0\leq D<e$, define $c_{\nu;a,b,D}$ by the following finite
multinomial rule. Sum over integers $n_{t,\alpha,d}\geq0$ satisfying
\[
 \begin{gathered}
 \sum_{\alpha,d}n_{t,\alpha,d}=\nu_t\quad\text{for each }t,\qquad
 \sum_{t,\alpha,d}d\,n_{t,\alpha,d}=D,\\
 \sum_{t,d,\alpha:\operatorname{Re}\alpha=1/2+\delta}
                                 n_{t,\alpha,d}=a,\qquad
 \sum_{t,d,\alpha:\operatorname{Im}\alpha=+\gamma}
                                 n_{t,\alpha,d}=b.
 \end{gathered}
\]
For these indices put
\[
 c_{\nu;a,b,D}=
 \sum_{\text{stated indices}}
   \left(\prod_{t=1}^{N-1}
             \frac{\nu_t!}{\prod_{\alpha,d}n_{t,\alpha,d}!}\right)
        \prod_{t,\alpha,d}f_{t;\alpha,d}^{\,n_{t,\alpha,d}}.
 \tag{OCS15}
\]
Empty index sets give zero and zero exponents give factors one.
Every original primary coefficient is present in this finite rule.
The ordinary multinomial expansion of OCS14 now proves
\[
 x(w)^\nu=\sum_{a,b=0}^k e^{\omega_{a,b}w}
                   \sum_{D=0}^{e-1}c_{\nu;a,b,D}w^D,
 \qquad
 \omega_{a,b}=(\beta_{a,b}-c)/i
                 =(2b-k)\gamma-i(2a-k)\delta.
 \tag{OCS16}
\]
Indeed $D\leq k(m-1)=e-1$ and the counts of each original
sign give exactly the displayed exponent. The sums with different
$(a,b)$ are distinct since their real and imaginary parts are
independently distinct.

\subsection{The actual finite observation matrix and its complete kernel}
Define the primary coefficient isomorphism
$\mathcal H:\mathbb C^q\to E$ by sending coordinate $(a,b,D)$
to $\varepsilon_{a,b}(S-\beta_{a,b})^D$, where
$\varepsilon_{a,b}$ is obtained from $\chi$ by the same truncated
Taylor-inverse construction as OCS12, now of order $e$.
The proof by the coprime factors is identical and shows that this
is a basis of all $q$ primary directions. Set
\[
 \mathfrak A_{\nu;(a,b,D)}=i^D D!\,c_{\nu;a,b,D},
 \qquad
 \mathfrak A:\mathbb C^q\longrightarrow
                            \mathbb C^{\mathcal I_{k,0}}.
 \tag{OCS17}
\]
Then, on the original spaces and with the original ordered target,
\[
 \boxed{\ L\mathcal H=\mathcal F_0\mathfrak A,\qquad
 K=\mathcal H(\ker\mathfrak A),\qquad
 \operatorname{rank}\Lambda=\operatorname{rank}\mathfrak A. }
 \tag{OCS18}
\]
Here is a proof giving every factorial and phase. In the primary
sum algebra,
\[
 e^{wY}[1]=\sum_{a,b}e^{\omega_{a,b}w}
          \sum_{D=0}^{e-1}\frac{i^{-D}w^D}{D!}
                  \varepsilon_{a,b}(S-\beta_{a,b})^D.
\]
The exact original tensor identity of OPG8 gives
$\mathcal V e^{wY}[1]=\mathbf v(w)^{\otimes k}$: the centre
$c=k/2$ splits into the $k$ original centres $1/2$, while
$U$ commutes with multiplication by $z_1$.
Applying the literal projector and then OCS9 and OCS16 gives
the right-hand curve with coefficients
$\mathcal F_0c_{a,b,D}$.
Exponential polynomials $w^De^{\omega_{a,b}w}$ with these distinct
exponents are linearly independent. To prove this directly,
apply $\prod_{(a',b')\ne(a,b)}(\partial_w-\omega_{a',b'})^e$
to any zero linear combination. It kills all the other exponents.
On the selected polynomial factor, each remaining operator is
$\partial_w+(\omega_{a,b}-\omega_{a',b'})$, invertible on the
polynomials of degree less than $e$ because its constant term
is nonzero and differentiation is nilpotent there. The surviving
polynomial must therefore vanish. Doing this for every pair proves
the asserted independence. Equating coefficients now gives
$L\mathcal H e_{a,b,D}=\mathcal F_0(i^DD!c_{a,b,D})$.
This proves the first equality in OCS18 on a basis. The other
two follow from injectivity of $\mathcal F_0$ and $\jmath$.

In particular the isomorphism from the actual sector image to
the original boundary is
\[
 \mathcal Q:\operatorname{im}\mathfrak A\xrightarrow{\ \sim\ }B,
 \qquad \mathcal Qx=\jmath^{-1}\mathcal F_0x,
 \qquad
 E/K\xrightarrow{\ \sim\ }\operatorname{im}\mathfrak A,
 \quad[z]\longmapsto\mathfrak A\mathcal H^{-1}z.
 \tag{OCS19}
\]
Both inverses exist on precisely their displayed images by OCS18.
These are full coefficient maps; the inherited Gram of
$\mathcal F_0$ is $J_0^*\mathcal GJ_0$ from OCS11.

For a completely explicit finite rank prescription, define
\[
 \Delta_r(\mathfrak A)=
   \sum_{\substack{I\subseteq\mathcal I_{k,0},\ J\subseteq\{(a,b,D)\}\\
                   |I|=|J|=r}}|\det\mathfrak A_{I,J}|^2,
 \qquad\Delta_0=1.
 \tag{OCS20}
\]
Every entry of every minor is the finite expression OCS12--17
in the actual periods and amplitude derivatives. Gaussian
elimination proves that a matrix has rank at least $r$ exactly
when one of its $r$ by $r$ minors is nonzero: pivot rows and
columns supply such a minor, and a nonzero minor supplies $r$
independent columns. Hence the actual rank is exactly
$\max\{r:\Delta_r(\mathfrak A)>0\}$. This evaluates its
finite deciding data without prescribing generic values or
asserting that any unexamined maximal minor is nonzero.

The Gamma vectors used in the existing endpoint increments have
the equally explicit coordinates
\[
 \jmath\lambda_d=\mathcal F_0\widehat\lambda_d,
 \qquad
 (\widehat\lambda_d)_\nu=
 \frac1{\sqrt{h_d^{(s)}}}\sum_{a,b=0}^k\sum_{D=0}^{e-1}
               c_{\nu;a,b,D}p_d^{(D)}(\omega_{a,b}),
 \quad h_d^{(s)}=(2\pi)^{s/2}d!(s/2)_d,
 \quad s\in\{1,k\}.
 \tag{OCS21}
\]
Taylor expansion of $p_d(Y)[1]$ in the sum primary basis gives
the coefficient $i^{-D}p_d^{(D)}(\omega_{a,b})/D!$.
Multiplication by OCS17 cancels exactly these factorials and
phases, proving OCS21. This also proves that the generating
formula of OPG9 becomes
$(1+z^2)^{-s/4}\mathcal F_0
 (x(\arctan z)^\nu)_{\nu\in\mathcal I_{k,0}}$,
with the original scalar exponent and both existing sources.

\subsection{Exact invariant dimensions and unconditional forced losses}
Put $d_{N,k,\tau}=|\mathcal I_{k,\tau}|$.
By OCS7--10 this is the dimension of the full symmetric sector,
and its exact finite character formula is
\[
 d_{N,k,0}=[T^k]\frac1N
       \sum_{h\mid N}\varphi(h)\frac{1-T}{(1-T^h)^{N/h}}.
 \tag{OCS22}
\]
To prove it, the trace generating series for $C^a$ on symmetric
tensors is
$\prod_{t=1}^{N-1}(1-T\zeta^{at})^{-1}$, by summing the
monomial basis OCS8. If $\zeta^a$ has order $h$, the list
$\zeta^{at}$ with $0\leq t<N$ contains every $h$th root
$N/h$ times. Their product is
$\prod_{t=0}^{N-1}(1-T\zeta^{at})=(1-T^h)^{N/h}$.
Removing the $t=0$ factor gives the series
$(1-T)/(1-T^h)^{N/h}$. There are exactly $\varphi(h)$
indices $a$ of that order, so averaging the traces of the
original finite projector proves OCS22. The inverse powers
give the same average by $a\mapsto-a$.
In ordinary finite binomial terms, for $k\geq1$ this is
\[
 d_{N,k,0}=\frac1N\sum_{h\mid N}\varphi(h)
 \left\{
 \mathbf1_{h\mid k}\binom{k/h+N/h-1}{N/h-1}
 -\mathbf1_{h\mid(k-1)}
                  \binom{(k-1)/h+N/h-1}{N/h-1}
 \right\}.
 \tag{OCS23}
\]
The binomial expansion of $(1-T^h)^{-N/h}$ follows by
counting weak compositions, proving the coefficient extraction
used here. It follows from the actual matrix OCS17 that
\[
 \operatorname{rank}\Lambda\leq\min(q,d_{N,k,0}),
 \qquad \dim K\geq\max(0,q-d_{N,k,0}).
 \tag{OCS24}
\]
These inequalities apply to the original map, since its full
image factorization was proved in OCS18; the ambient dimension
has not been equated with its actual rank.
There is also an unconditional nonzero observation for the
original $k\geq5$. Each entire function $x_t(w)$ in OCS14 is
nonzero. Otherwise its nonzero coordinate covector would
annihilate every derivative at zero of
$\Pi U\exp(w(A_h-1/2)/i)[1]$. Those derivatives span the full
space, since $[1]$ is cyclic for multiplication by $z_1$ and
$\Pi U$ is invertible. A product of finitely many nonzero entire
functions is nonzero. The zero sector is nonempty: take the
weights $1,1,N-2$ and then $(k-3)/2$ pairs $1,N-1$.
Their sum is a multiple of $N$. Its monomial curve in OCS9
is nonzero, so $\operatorname{rank}\Lambda\geq1$.

For $m=1$, $N=5$, and every $k\geq0$, write
\[
 D_k=\binom{k+3}{3},\qquad
 \epsilon_k=\begin{cases}1&k\equiv0\pmod5,\\-1&k\equiv1\pmod5,\\
                         0&\text{otherwise}.
             \end{cases}
 \qquad
 d_{5,k,0}=\frac{D_k+4\epsilon_k}{5},\quad
 d_{5,k,\tau}=\frac{D_k-\epsilon_k}{5}\quad(\tau\ne0).
 \tag{OCS25}
\]
For each of the four nonidentity powers of $C$, the trace
series is $(1-T)/(1-T^5)$, whose coefficient is $\epsilon_k$.
The weight-$\tau$ projector introduces the Fourier multiplier;
the sum of its four nontrivial fifth roots is $4$ when
$\tau=0$ and $-1$ otherwise. This proves both formulas.
For the requested original orders the exact counts and resulting
kernel lower bounds are
\[
\begin{array}{c|r|r|r|r}
 k&D_k&d_{5,k,0}&q=(k+1)^2&\max(0,q-d_{5,k,0})\\\hline
 5&56&12&36&24\\
 9&220&44&100&56\\
 13&560&112&196&84\\
 17&1140&228&324&96\\
 21&2024&404&484&80\\
 25&3276&656&676&20\\
 29&4960&992&900&0
\end{array}\tag{OCS26}
\]
These are integer evaluations of OCS25. The strict inequality
$d_{5,k,0}<q$ holds for $k=4l+1$, $l\geq1$, exactly for the
first six orders in the table. They exhaust $k\leq25$ in this
progression. For $k\geq29$,
\[
 D_k-5(k+1)^2=\frac{(k+1)(k^2-25k-24)}6\geq460>4,
\]
where $k^2-25k-24$ is positive and increasing from $k=29$.
Even the most negative correction $4\epsilon_k\geq-4$
therefore gives $d_{5,k,0}>q$. This proves the threshold
without extrapolating the finite table.

\subsection{The actual simple-quartet quadric and a stronger rank bound}
Continue with $m=1$. Order the four roots as
$(++),(+-),(-+),(--)$ and let
$\mathcal E_{\rm prim}:\mathbb C^4\to E_h$ be the primary
basis isomorphism whose columns are the four $[E_{\alpha,0}]$.
Define the actual invertible matrix and the original exponential
quartet by
\[
 \mathsf H=F^{-1}\Pi U\mathcal E_{\rm prim},\qquad
 t_{\epsilon,\eta}(w)=e^{(\eta\gamma-i\epsilon\delta)w},
 \qquad x(w)=\mathsf Ht(w).
 \tag{OCS27}
\]
Every column of $\mathsf H$ is OCS13 with $d=0$, namely the
Fourier transform of its original primary period times its
actual nonzero $v_h(\alpha)$. Each displayed map is invertible,
so $\mathsf H$ is invertible without a genericity assumption.
Put $\mathcal R_k=\mathbb C[x_1,x_2,x_3,x_4]_k$, the homogeneous
polynomials of degree $k$, and define
\[
 Q_0(t)=t_{++}t_{--}-t_{+-}t_{-+},\qquad
 Q(x)=Q_0(\mathsf H^{-1}x).
 \tag{OCS28}
\]
This is the literal transformed quadratic relation of the
original curve; its coefficients are determined by the periods
and unit through OCS27. In particular $Q\ne0$.

The kernel of homogeneous restriction to that curve is exactly
\[
 \ker\{\mathcal R_k\to\operatorname{Hol}(\mathbb C),
                         f\mapsto f(x(w))\}
            =Q\mathcal R_{k-2}\quad(k\geq2).
 \tag{OCS29}
\]
Here is a complete verification. First work in the $t$ coordinates.
A degree-$k$ monomial has exponent
$(2b-k)\gamma-i(2a-k)\delta$, where $a$ is its count of
first sign $+$ and $b$ its count of second sign $+$.
These $(k+1)^2$ complex exponents are distinct: equality of
their real parts forces $b$ equal and equality of their
imaginary parts forces $a$ equal. Equivalently, the two
frequencies $2\gamma$ and $-2i\delta$ admit no nonzero
integer relation, by their real and imaginary parts. Thus
the exponentials have no extra polynomial relation of this
homogeneous type; their linear independence follows, for
example, by differentiating $0,\ldots,(k+1)^2-1$ times and
using the Vandermonde determinant of the distinct exponents.
Every pair $(a,b)$ occurs: the count $r_{++}$ may be any integer
between $\max(0,a+b-k)$ and $\min(a,b)$, a nonempty interval,
and the other three counts are $a-r_{++}$, $b-r_{++}$,
and $k-a-b+r_{++}$.

Within any fixed pair $(a,b)$, two monomial exponent vectors
differ by an integer multiple of $(1,-1,-1,1)$. After their
common monomial is factored out, their difference is therefore
the difference of equal powers of $t_{++}t_{--}$ and
$t_{+-}t_{-+}$. The identity
$A^r-B^r=(A-B)\sum_{j=0}^{r-1}A^{r-1-j}B^j$ shows that
this difference is divisible by $Q_0$. Any polynomial
vanishing on the curve has coefficient sum zero within
each pair $(a,b)$, by the proved exponential independence.
It is thus a sum of these monomial differences and is
divisible by $Q_0$. Its quotient is homogeneous of degree
$k-2$. Conversely $Q_0(t(w))=0$, since each of its two
products is one. This proves the assertion in $t$
coordinates, and the invertible substitution OCS27 proves
OCS29. In particular
$\dim(\mathcal R_k/Q\mathcal R_{k-2})=(k+1)^2=q$;
no information about the original curve has been replaced
by a freely chosen quadric.

Let $\mathcal I_k\subseteq\mathcal R_k$ be the span of the
weight-zero monomials $x^\nu$, using the weights $1,2,3,4$
from the exact Fourier basis. The coordinate functions of
the original projected curve in OCS9 are precisely these
monomials restricted to $x(w)$. The span of the vector
curve $Le^{wY}[1]$ equals $\operatorname{im}L$: its
derivatives at zero contain $LY^d[1]$, $0\leq d<q$, and
$Y^d[1]$ is a basis because $Y$ is an invertible affine
change of multiplication by $S$ in the degree filtration.
The rank of that vector curve equals the dimension of
the span of its coordinate functions, by the equality of
row and column ranks of its finite exponential coefficient
matrix. Hence OCS29 proves the exact equality
\[
 \operatorname{rank}\Lambda
   =d_{5,k,0}-\dim(\mathcal I_k\cap Q\mathcal R_{k-2}).
 \tag{OCS30}
\]

Fix any multiplicative monomial order, for instance
lexicographic order, and write $x^\mu$ for the leading
monomial of the actual $Q$, with $|\mu|=2$. Choose a
row-echelon basis of the finite-dimensional intersection
in OCS30 in descending monomial order. Its leading
monomials are distinct, all have weight zero, and all
are divisible by $x^\mu$: the leading monomial of
$Qf$ is the product of those of $Q$ and $f$ because
the order is multiplicative and the coefficient field
has no zero divisors. Thus each such leading monomial
is $x^\mu x^\eta$ for a distinct degree-$k-2$ monomial
of weight $-\operatorname{wt}(\mu)$. Counting this set
and using OCS30 gives
\[
 \boxed{\quad
 d_{5,k,0}-d_{5,k-2,-\operatorname{wt}(\mu)}
 \ \leq\ \operatorname{rank}\Lambda
 \ \leq\ \min(q,d_{5,k,0}).\quad}
 \tag{OCS31}
\]
The left side depends only on the actual leading monomial
of the actual period quadric, and remains valid whatever
its other coefficients are. A uniform consequence retaining
the strongest count independent of that monomial is
\[
 r_k^-:=d_{5,k,0}-\max_{\tau\in\mathbb Z/5}d_{5,k-2,\tau}
       =\left\lfloor\frac{(k+1)^2}{5}\right\rfloor
                                  +\mathbf1_{5\mid k}
       \leq\operatorname{rank}\Lambda.
 \tag{OCS32}
\]
For the equality, $D_k-D_{k-2}=(k+1)^2$, as direct
subtraction of the two cubic binomials shows. Formula
OCS25 then gives
\[
 r_k^-=\frac{(k+1)^2+4\epsilon_k
                   -\max(4\epsilon_{k-2},-\epsilon_{k-2})}{5}.
\]
For $k$ congruent respectively to $0,1,2,3,4$ modulo five,
its numerator correction is respectively $4,-4,-4,-1,0$.
The corresponding residues of $(k+1)^2$ are $1,4,4,1,0$.
These five cases prove OCS32 for every $k\geq2$.
For the original orders in OCS26 the complete unconditional
intervals are therefore
\[
\begin{array}{c|r|r|r|r}
 k&r_k^-&\min(q,d_{5,k,0})&\max(0,q-d_{5,k,0})&q-r_k^-\\\hline
 5&8&12&24&28\\
 9&20&44&56&80\\
 13&39&112&84&157\\
 17&64&228&96&260\\
 21&96&404&80&388\\
 25&136&656&20&540\\
 29&180&900&0&720
\end{array}\tag{OCS33}
\]
The middle two rank columns bound $\operatorname{rank}\Lambda$;
the last two columns bound $\dim K$ from below and above,
respectively. No endpoint of either interval has been assigned
as the actual value without its original minors.

\subsection{A second exact finite matrix for the simple-quartet kernel}
The quadric calculation also gives the actual kernel a
polynomial quotient model. Write $Q(x)=\sum_{|\mu|=2}q_\mu x^\mu$.
Let $T_Q$ be the full multiplication matrix
$\mathcal R_{k-2}\to\mathcal R_k$ in the monomial bases:
\[
 (T_Q)_{\nu,\eta}=q_{\nu-\eta},
 \quad |\nu|=k,\quad|\eta|=k-2,
 \qquad q_{\nu-\eta}=0\text{ if an entry of }\nu-\eta<0.
\]
Let $T_{\rm out}$ be its rows of nonzero cyclic weight.
Every coefficient $q_\mu$ is explicitly fixed by
$Q_0(\mathsf H^{-1}x)$, and the inverse matrix may be
computed as $\operatorname{adj}\mathsf H/\det\mathsf H$,
with nonzero denominator proved in OCS27. Multiplication
by $Q$ is injective. The condition that $Qf$ belongs
to $\mathcal I_k$ is exactly $T_{\rm out}f=0$.
Thus OCS30 gives the further exact rank identities
\[
 \begin{split}
 \operatorname{rank}\Lambda
   &=d_{5,k,0}-D_{k-2}+\operatorname{rank}T_{\rm out},\\
 \dim K&=D_k-d_{5,k,0}-\operatorname{rank}T_{\rm out}.
 \end{split}\tag{OCS34}
\]
The ranks of this concrete matrix are decided by its finite
minors, by the same elimination proof as OCS20. This is
an alternative exact calculation from the actual periods,
not an assumption about the quadric's coefficients.

To give the typed map behind the kernel equality, let
$E^*$ denote the complex linear dual. Define
$\Theta:\mathcal R_k\to E^*$ as follows: $\Theta(f)$ is
the unique covector such that
\[
 \Theta(f)(e^{wY}[1])=f(x(w))\quad\text{for all }w.
 \tag{OCS35}
\]
For $m=1$ every function on the right is a linear
combination of the $q$ distinct exponentials in OCS16.
The primary expansion of $e^{wY}[1]$ gives exactly those
exponentials as a basis of covector functions. This proves
existence and uniqueness. The proof of OCS29 proves
surjectivity and $\ker\Theta=Q\mathcal R_{k-2}$.
Moreover $\Theta(\mathcal I_k)=\operatorname{im}L^\vee$,
where $L^\vee$ denotes the complex linear dual map:
the monomial coefficient functionals in the basis OCS8
pull back along $L$ to exactly OCS35 for $x^\nu$ of
weight zero. These functionals span the full dual of
the actual image, since every linear functional on a
subspace of a finite-dimensional space extends to the
whole space. Restriction of $E^*$ to $K$ has kernel
$\operatorname{im}L^\vee$: a covector vanishing on $K$
factors uniquely through $E/K\simeq\operatorname{im}L$.
Thus the exact maps give
\[
 \boxed{\quad
 \operatorname{coker}T_{\rm out}
 \simeq\mathcal R_k/(\mathcal I_k+Q\mathcal R_{k-2})
 \xrightarrow{\ \Theta|_K\ }K^*.
 \quad}\tag{OCS36}
\]
For the first isomorphism, use the direct monomial splitting
of $\mathcal R_k$ into its zero and nonzero sectors; after
quotienting by $\mathcal I_k$, multiplication by $Q$
becomes precisely $T_{\rm out}$. For the second, the
surjectivity and kernel were just proved. This identifies
the actual lost directions by a full quotient map.

\subsection{All primary action directions remain in the mixed calculation}
For arbitrary original $m\geq1$, retain the full primary
matrix $\mathfrak A$ in OCS17. In the basis $\mathcal H$
the original sum action is exactly
\[
 J e_{a,b,D}=\omega_{a,b}e_{a,b,D}
       +\begin{cases}i^{-1}e_{a,b,D+1}&D<e-1,\\0&D=e-1,
         \end{cases}
 \qquad Y\mathcal H=\mathcal HJ.
 \tag{OCS37}
\]
This follows by multiplying
$\varepsilon_{a,b}(S-\beta_{a,b})^D$ by
$(S-c)/i=\omega_{a,b}+(S-\beta_{a,b})/i$ in its original
primary algebra. Consequently the actual possible action
leaving the cyclic kernel is the fully specified map
\[
 \ker\mathfrak A\longrightarrow\operatorname{im}\mathfrak A,
 \qquad x\longmapsto\mathfrak A Jx,
 \qquad
 \Lambda Y\mathcal Hx=\mathcal Q\mathfrak A Jx.
 \tag{OCS38}
\]
The codomain assertion follows from $\mathfrak A Jx$ being
a value of the same matrix $\mathfrak A$ on $Jx$.
For the original inclusion $I:K\hookrightarrow E$, this
is the OPG22 defect transported through the exact kernel
isomorphism in OCS18. In particular none of the discarded
sectors has been declared invariant under $Y$.

The part of the kernel undetected by all original iterates
is exactly computed by the stacked matrix
\[
 \mathfrak O=\begin{bmatrix}
          \mathfrak A\\\mathfrak AJ\\\vdots\\\mathfrak AJ^{q-1}
          \end{bmatrix},\qquad
 \mathcal K_\infty=\mathcal H\ker\mathfrak O,
 \quad
 \dim\mathcal K_\infty=q-\operatorname{rank}\mathfrak O,
 \quad
 \dim(K/\mathcal K_\infty)=
        \operatorname{rank}\mathfrak O-\operatorname{rank}\mathfrak A.
 \tag{OCS39}
\]
Indeed $\mathfrak A J^a x$ is the exact sector-coordinate
value of $\Lambda Y^a\mathcal Hx$ by OCS18 and OCS37.
The polynomial $\chi(c+iT)$ annihilates $J$ and has
degree $q$ with leading coefficient $i^q$. It expresses
every higher power in the span of the first $q$, so the
stacked kernel is precisely the intersection over all
iterates and is $J$-invariant. Rank-nullity proves both
dimension formulas. Its quotient identifies the directions
that the original one-step observation loses and subsequent
original iterates detect.

Equations OCS17--21, OCS34--39 supply actual finite maps
and deciding minors for the original observation and its
kernel. The forced kernel losses in OCS26 and the rank
intervals OCS33 concern this unchanged projector and these
unchanged periods. The endpoint volume calculation still
uses the full positive sums of the actual vectors
$\lambda_d$ in OPG16--18 and their exact inherited
coordinates OCS11 and OCS19. No volume growth estimate
is deduced from these dimension bounds alone.

\end{document}
